\section{The idea}
\label{sec:idea}
Teleoperation is limited by robot time: every teleoperated episode occupies the robot as well as the operator. A glove demonstration occupies only the collector. Ours is at least as fast per demonstration as teleoperation of the same task (2.0--3.1 against 1.6 per minute, \S\ref{sec:system}), and none of that time is the robot's: collectors can work in parallel, each one needs a \$22 glove and a \$122 camera, and a new workspace needs only four printed tiles taped around it.

Every UMI-style device answers two questions: what does the policy see, and where was the gripper. The standard answer to the first is a camera on the wrist, which by construction sees what the robot's wrist camera will see. The standard answer to the second is a pose tracker on the same wrist, and every tracker has a price. SLAM over the wrist camera \citep{chi2024umi,campos2021orbslam3} needs a mapping pass for every new scene and enough texture in it, and loses track under occlusion and fast motion; an off-the-shelf visual-inertial tracker adds a second device, a RealSense T265, to the handheld gripper \citep{fastumi2024}; a VR controller in or on the gripper \citep{yubi2026,tinyumi2026} adds weight to the hand (\SI{119}{g} of YUBI's \SI{319}{g}) and brings a tracking system of its own, whose accuracy has to be measured separately \citep{vrtracking2024,quest2rehab2022}; a fixed 4K camera over the workspace, the choice of our earlier glove \citep{yamumi2026}, gives absolute labels but ties collection to the camera's mount, and the collector's own body hides the hand from it, \SI{8}{}--\SI{22}{\percent} of frames per session in the data below (\S\ref{sec:labels}). The wrist camera has prices of its own. It puts electronics on the glove: a camera, a cable and a stream to synchronise with the tracker. And it frames the jaws and little else, so the goal of a carry, here the bin, can leave the policy's view while the task still depends on it: a hand-centric view trades observability of the scene for generalisation \citep{hsu2022hands}.

Two of the open problems this workshop names \citep{umiarena2026} are in-the-wild tracking that does not rest on SLAM, and the viewpoint gap between what a handheld device records and what the robot sees. Our answer to both is one camera, worn. The collector's head tracks the hand by habit, since gaze goes to each object just before the hand acts on it \citep{land2001eye}; it keeps the gripper roughly centred at arm's length and is never occluded by the torso. From its view the fiducial ball gives the pose, the jaw tags give the width, and printed AprilTag \citep{olson2011apriltag} tiles taped on the table give the world frame, so the glove itself carries no electronics at all, and the scene around the gripper, bin included, stays in frame. The same camera is the policy's observation. On the robot a small servo-driven printed arm carries it and imitates the head from the robot's joint state alone, identically during teleoperation and rollouts, so training and deployment share one viewpoint distribution. The price is a moving camera at deployment; the claim is that this price buys the wrist camera, the tracker and the cable, and that the policy is as good. A fixed camera was kept only as an ablation.

Two lines of work are closest. EgoMimic \citep{kareer2025egomimic} co-trains one policy on robot data and egocentric human video from head-worn glasses, with the hand's pose tracked from the video; our collector holds a replica of the robot's gripper, so the contact is the robot's and the jaw width is measured rather than estimated. Active-vision teleoperation \citep{chuang2025avaloha,cheng2024opentelevision} moves a robot camera with the operator's head through a VR headset; our positioner has no operator to follow during rollouts, so it replays a model of how the collector's head moved, from the robot's joint state alone, identically in teleoperation and deployment.
